% ============================================================
%  Week 6 — Propensity Score Methods
%  Compile standalone:  pdflatex chapter6_new.tex
%  Compile via master:  pdflatex causal_inference_master.tex
% ============================================================
\documentclass[master.tex]{subfiles}

\begin{document}

% ------------------------------------------------------------
%  Title block (standalone) / chapter heading (master)
% ------------------------------------------------------------
\ifSubfilesClassLoaded{%
  \begin{center}
    {\LARGE\bfseries\color{isublue}
     Chapter 6: Propensity Score Methods}\\[6pt]
    {\large Theory and Methods in Causal Inference}\\[4pt]
    {\normalsize Department of Statistics, Iowa State University}
  \end{center}
  \bigskip
  \hrule height 1.2pt \relax
  \smallskip
}{%
  \chapter{Propensity Score Methods}
}

% ------------------------------------------------------------
%  Learning objectives
% ------------------------------------------------------------
\begin{tcolorbox}[defstyle, title={Learning Objectives}]
By the end of this lecture, students should be able to:
\begin{enumerate}
  \item Define the propensity score $\pi(X) = P(T{=}1 \mid X)$ and
    explain why it is a \emph{balancing score}.
  \item State and prove the Balancing Property ($T \indep X \mid \pi(X)$)
    and the Propensity Score Theorem (strong ignorability given $X$
    implies strong ignorability given $\pi(X)$).
  \item Describe standard methods for estimating the propensity score
    (logistic regression, machine learning) and the practical pitfalls
    of each.
  \item Derive the IPW identification formula for the ATE and explain
    the role of the Horvitz--Thompson representation.
  \item Distinguish the IPW estimator (targets ATE) from the
    propensity-score matching estimator (targets ATT).
  \item Define strict overlap, describe the practical consequences of
    near-violations, and apply trimming strategies.
  \item Articulate the fundamental limitation of propensity-score
    methods --- the unconfoundedness assumption --- and explain why this
    motivates instrumental variables (Chapter~7).
\end{enumerate}
\end{tcolorbox}


% ============================================================
\section{Motivation: The Curse of Dimensionality}
% ============================================================

Chapter~5 established that under strong ignorability, the ATE is
identified by the back-door adjustment formula, and developed three
estimation strategies: regression adjustment, standardisation, and
stratification.  All three require, in practice, that we condition on a
covariate vector $X$.  When $X$ is low-dimensional, this is feasible.
When $X$ has many components, it is not.

\medskip\noindent\textbf{The problem.}  Stratification requires cells
$\{X = x\}$ to contain both treated and control units.  With $p$ binary
covariates, there are $2^p$ cells.  With $p = 10$, that is 1024 cells
--- far more than most datasets can support.  Regression adjustment
avoids the cell-count problem but requires a correctly specified model
for $\E[Y \mid T, X]$, which becomes harder to specify reliably as $p$
grows.

\medskip\noindent\textbf{The solution.}  Rosenbaum \& Rubin (1983)
showed that it is sufficient to condition on a single scalar function
of $X$: the \emph{propensity score} $\pi(X) = P(T{=}1 \mid X)$.  This
reduces a $p$-dimensional adjustment problem to a one-dimensional one,
without any loss of identification power.

% ============================================================
\section{The Propensity Score and Its Balancing Properties}
% ============================================================

The central insight of \citet{rosenbaum1983central} is that the
$p$-dimensional adjustment problem of Chapter~5 reduces to a
one-dimensional one.  The right abstraction is the
\emph{balancing score}: any function of $X$ that renders treatment
assignment conditionally independent of the covariates.  The
propensity score is the canonical---and coarsest---such function.

\subsection{Balancing Scores and the Propensity Score}

\begin{defn}{Balancing Score \citep{rosenbaum1983central}}
\label{w6:defn:balancing-score}
A function $b(X)$ is called a \textbf{balancing score} if
\[
  T \indep X \mid b(X).
\]
\end{defn}

\noindent The full covariate vector $X$ is trivially a balancing score:
conditioning on $X$ renders treatment independent of itself.  The
propensity score is the \emph{coarsest} balancing score
(Lemma~\ref{w6:lem:coarsest} below), and therefore provides the
greatest dimension reduction without sacrificing identification.

\begin{defn}{Propensity Score \citep{rosenbaum1983central}}
\label{w6:defn:ps}
For a binary treatment $T \in \{0,1\}$ and observed covariates $X$,
the \textbf{propensity score} is
\[
  \pi(X) = P(T{=}1 \mid X).
\]
\end{defn}

\noindent In a randomized experiment, $\pi(X) = 0.5$ for all units.  In
an observational study, $\pi(X)$ varies with $X$ and must be estimated
from data.

\medskip\noindent\textbf{Connection to strong ignorability.}
More explicitly, the propensity score can be written as
$P\{T{=}1 \mid X,\, Y(1),\, Y(0)\}$, the probability of treatment
given \emph{all} available information including the potential outcomes.
Under strong ignorability $(Y(0),Y(1)) \indep T \mid X$, the potential
outcomes carry no additional information about assignment beyond $X$,
so
\[
  P\bigl\{T{=}1 \mid X,\, Y(1),\, Y(0)\bigr\}
  = P\{T{=}1 \mid X\}
  = \pi(X).
\]
The gap between the two expressions is precisely what unconfoundedness
rules out.

\begin{thm}{The Propensity Score is a Balancing Score \citep{rosenbaum1983central}}
\label{w6:thm:balance}
$T \indep X \mid \pi(X)$.
\end{thm}

\begin{proof}
It suffices to show $P(T{=}1 \mid X, \pi(X)) = P(T{=}1 \mid \pi(X))$.
Since $\pi(X)$ is a deterministic function of $X$, conditioning on both
$X$ and $\pi(X)$ is the same as conditioning on $X$ alone:
\[
  P\bigl(T{=}1 \mid X, \pi(X)\bigr) = P(T{=}1 \mid X) = \pi(X)
  = P\bigl(T{=}1 \mid \pi(X)\bigr). \qquad\square
\]
\end{proof}


\begin{remark}[Covariate balance and design-stage analysis]
\label{w6:rem:balance-check}
Theorem~\ref{w6:thm:balance} has a practical consequence that
\citet{imbens2015causal} call \emph{design before analysis}: covariate
balance can be assessed before the outcome data are consulted.
The balancing property implies
\[
  f\bigl(X \mid T{=}1,\, \pi(x)\bigr) = f\bigl(X \mid T{=}0,\, \pi(x)\bigr).
\]
Stratifying on a discretised $\hat\pi(X)$ and comparing covariate
distributions within strata---via standardized mean differences or
$t$-tests---diagnoses whether the propensity score model achieves
adequate balance, without ever looking at $Y$.  Because the diagnostic
does not involve the outcome, it cannot be used to ``fish'' for a
preferred model specification, making the analysis more objective.
\end{remark}

\begin{remark}
\label{w6:rem:unobserved}
Theorem~\ref{w6:thm:balance} is a statement about the \emph{observed}
covariate distribution.  It says that within strata of $\pi(X)$, the
distribution of $X$ is the same in the treated and control groups.  It
does \emph{not} say anything about unobserved confounders $U$: if $U$
affects both $T$ and $Y$ and is not captured by $X$, the balancing
property fails to close the back-door path through $U$.
\end{remark}

\subsection{The Propensity Score Theorem}

The balancing property has a fundamental identification consequence:
unconfoundedness at $X$ implies unconfoundedness at \emph{any}
balancing score of $X$, and in particular at the scalar $\pi(X)$.

\begin{thm}{Propensity Score Theorem \citep{rosenbaum1983central}}
\label{w6:thm:ps}
Suppose $(Y(0), Y(1)) \indep T \mid X$.  Then for any balancing score
$b(X)$,
\[
  \bigl(Y(0),\, Y(1)\bigr) \;\indep\; T \;\Big|\; b(X).
\]
In particular, taking $b(X) = \pi(X)$: strong ignorability given $X$
implies strong ignorability given $\pi(X)$ alone, provided
$0 < \pi(X) < 1$ almost surely.
\end{thm}

\begin{proof}
We prove the result for $b(X) = \pi(X)$; the general case follows
by noting that $\pi(X) = g(b(X))$ for some measurable function $g$
(Lemma~\ref{w6:lem:coarsest}), so the same argument applies
with $\pi(X)$ replaced by $b(X)$ throughout.

By the law of iterated expectations, for any $t \in \{0,1\}$:
\begin{align*}
  P\bigl(T{=}1 \mid Y(t),\, \pi(X)\bigr)
  &= \E\bigl[P(T{=}1 \mid Y(t), X) \mid Y(t), \pi(X)\bigr] \\
  &= \E\bigl[P(T{=}1 \mid X) \mid Y(t), \pi(X)\bigr]
     \quad\text{(ignorability: }T \indep Y(t) \mid X\text{)} \\
  &= \E\bigl[\pi(X) \mid Y(t), \pi(X)\bigr] = \pi(X).
\end{align*}
Since $P(T{=}1 \mid Y(t), \pi(X)) = \pi(X) = P(T{=}1 \mid \pi(X))$, we
have $T \indep Y(t) \mid \pi(X)$.
\end{proof}

\begin{remark}
\label{w6:rem:conditional-exp}
Theorem~\ref{w6:thm:ps} gives an alternative identification formula:
\[
  \E[Y(t)] = \E\bigl[\E[Y \mid \pi(X),\, T{=}t]\bigr].
\]
A common misconception is that this means the two conditional
regressions $\E[Y \mid \pi(X), T{=}t]$ and $\E[Y \mid X, T{=}t]$
are equal.  They are not: they are different functions and their
pointwise values generally differ.  The theorem guarantees only that
their marginal expectations agree:
\[
  \E\bigl[\E[Y \mid \pi(X),\, T{=}t]\bigr]
  = \E\bigl[\E[Y \mid X,\, T{=}t]\bigr]
  = \E[Y(t)].
\]
In particular, regressing $Y$ on $\pi(X)$ and $T$ is \emph{not}
equivalent to regressing $Y$ on $X$ and $T$; both identify the ATE,
but through different nuisance functions.
\end{remark}

\begin{lemma}{The Propensity Score is the Coarsest Balancing Score \citep{imbens2015causal}}
\label{w6:lem:coarsest}
If $b(X)$ is any balancing score---that is, $T \indep X \mid b(X)$---then
$\pi(X)$ is a function of $b(X)$: $\pi(X) = g(b(X))$ for some measurable
function $g$.
\end{lemma}

\begin{proof}
Since $b(X)$ is a balancing score, $T \indep X \mid b(X)$, which means
$P(T{=}1 \mid X,\, b(X)) = P(T{=}1 \mid b(X))$.  But $b(X)$ is a
function of $X$, so conditioning on $b(X)$ is already implicit in
conditioning on $X$:
\[
  \pi(X) = P(T{=}1 \mid X) = P\bigl(T{=}1 \mid X,\, b(X)\bigr)
        = P\bigl(T{=}1 \mid b(X)\bigr).
\]
Hence $\pi(X) = P(T{=}1 \mid b(X))$, which is a measurable function of
$b(X)$ alone.
\end{proof}

\noindent Lemma~\ref{w6:lem:coarsest} says that $\pi(X)$ is the
\emph{smallest} sufficient statistic for $T$ among all functions of
$X$: any other balancing score retains at least as much information
about $X$ as the propensity score does, making $\pi(X)$ the optimal
choice for dimension reduction.

\begin{tcolorbox}[defstyle, title={The Dimension Reduction}]
Theorem~\ref{w6:thm:ps} says that instead of adjusting for the full
high-dimensional vector $X$, it is sufficient to adjust for the scalar
$\pi(X)$.  This reduces the \emph{curse of dimensionality} in covariate
adjustment from a $\dim(X)$-dimensional problem to a one-dimensional
one.  Identification of the ATE follows immediately:
\[
  \tau_{\mathrm{ATE}}
  = \E\bigl[\E[Y \mid T{=}1, \pi(X)] - \E[Y \mid T{=}0, \pi(X)]\bigr].
\]
\end{tcolorbox}

% ============================================================
\section{Estimation of the Propensity Score}
% ============================================================

The true propensity score $\pi(X) = P(T{=}1 \mid X)$ is \emph{unknown}
in observational studies and must be estimated from the data.

\subsection*{Logistic Regression}

The classical approach models:
\[
  \log\frac{\pi(X)}{1 - \pi(X)} = X^\top \beta,
\]
and estimates $\beta$ by maximum likelihood.  The fitted values
$\hat\pi(X_i) = \sigma(X_i^\top \hat\beta)$, where $\sigma$ is the
logistic function, are used in place of the true propensity score.

\subsection*{Machine Learning Estimators}

When $X$ is high-dimensional or when the true propensity score is
non-linear, machine learning methods --- gradient boosted trees, random
forests, regularised logistic regression (LASSO, ridge) --- can estimate
$\pi(X)$ more flexibly.  A critical issue is \emph{overfitting}: an
estimator that perfectly separates treated and control units in-sample
will produce estimated propensity scores near 0 or 1 everywhere,
destroying overlap.  Cross-fitting (Chapter~11) addresses this by
estimating nuisance functions on held-out folds.

\begin{warn}{Propensity Score Estimation is a Nuisance, Not the Goal}
The propensity score is estimated in order to construct a good estimator
of $\tau_{\mathrm{ATE}}$.  Balancing tests --- checking whether
$\E[X \mid T{=}1, \hat\pi(X) \approx p]
 \approx \E[X \mid T{=}0, \hat\pi(X) \approx p]$ ---
diagnose whether the estimated score has achieved covariate balance, but
they do \emph{not} test whether unobserved confounders are balanced.
\end{warn}

% ============================================================
\section{Matching on the Propensity Score}
% ============================================================

The propensity score also motivates \emph{matching}: for each treated
unit $i$, find a control unit $j$ with $\hat\pi(X_j) \approx
\hat\pi(X_i)$, and use $Y_j$ as a proxy for $Y_i(0)$
\citep{abadie2006large, abadie2016matching}.  The resulting
estimator of the ATT is:
\[
  \hat\tau_{\mathrm{ATT}}
  = \frac{1}{n_1}\sum_{i:\, T_i=1}
    \Bigl[Y_i - Y_{\hat\jmath(i)}\Bigr],
\]
where $\hat\jmath(i)$ is the index of the matched control unit.

\subsection{One-to-One Nearest Neighbour Matching}

The most common matching algorithm is \emph{nearest-neighbor} (1:1)
matching \citep{abadie2006large}: for each treated unit $i$, find:
\[
  \hat\jmath(i) = \arg\min_{j:\, T_j=0}
    \bigl|\hat\pi(X_i) - \hat\pi(X_j)\bigr|.
\]
Matched control units may be selected with or without replacement.
Matching with replacement reduces bias (each control unit is available
to every treated unit) but increases variance (a small number of
control units may be used many times).

\subsection{Caliper Matching}

Caliper matching restricts matches to pairs within a maximum distance
$\delta$ (the ``caliper'') on the propensity score scale:
\[
  \hat\jmath(i) = \arg\min_{j:\, T_j=0,\;
    |\hat\pi(X_i) - \hat\pi(X_j)| \le \delta}
    \bigl|\hat\pi(X_i) - \hat\pi(X_j)\bigr|.
\]
Treated units with no control match within the caliper are discarded.
This restricts the estimand to the subset of treated units with good
covariate overlap.

\begin{remark}
The key idea behind propensity score matching is that it implicitly fits
a \emph{nonparametric regression} of $Y$ on $\pi(X)$ separately within
each treatment arm.  By Theorem~\ref{w6:thm:ps}, $(Y(0), Y(1)) \indep T
\mid \pi(X)$, so the conditional mean $\E[Y(t) \mid \pi(X) = p]$ is
identified from observed data for each $t \in \{0, 1\}$.  Nearest-neighbor
matching approximates this conditional mean at each observed $\hat\pi(X_i)$
by averaging the outcomes of close matches, effectively performing local
constant regression on the propensity score.  The dimension reduction from
Theorem~\ref{w6:thm:ps} is therefore what makes nonparametric estimation
feasible: instead of smoothing over the full covariate vector $X$, one
smooths over the scalar $\pi(X)$.
\end{remark}

\begin{remark}
Propensity score matching provides \emph{design transparency}: by
inspecting matched pairs, the analyst can directly verify whether
treated and control units are similar on observed covariates.
Comparing $\bar X_1$ and $\bar X_0$ within matched pairs via
standardized mean differences is a standard pre-analysis diagnostic
that can be conducted without consulting the outcome $Y$.
\end{remark}

% ============================================================
\section{Inverse Probability Weighting}
\label{w6:sec:ipw}
% ============================================================

\subsection{The IPW Identification Formula}

Under strong ignorability, the ATE admits an alternative representation
as an \emph{inverse probability weighting} (IPW) formula:
\[
  \tau_{\mathrm{ATE}}
  = \E\!\left[\frac{T \cdot Y}{\pi(X)}\right]
  - \E\!\left[\frac{(1-T)\cdot Y}{1-\pi(X)}\right].
\]

\begin{thm}{IPW Identification}
\label{w6:thm:ipw}
Under strong ignorability ($(Y(0), Y(1)) \indep T \mid X$ and
$0 < \pi(X) < 1$):
\[
  \E\!\left[\frac{T \cdot Y}{\pi(X)}\right] = \E[Y(1)].
\]
\end{thm}

\begin{proof}
\[
  \E\!\left[\frac{T Y}{\pi(X)}\right]
  = \E\!\left[\frac{T Y(1)}{\pi(X)}\right]
  = \E\!\left[\E\!\left[\frac{T Y(1)}{\pi(X)} \;\Big|\; X\right]\right]
  = \E\!\left[\frac{Y(1)}{\pi(X)}\, \E[T \mid X]\right]
  = \E\!\left[\frac{Y(1)}{\pi(X)}\, \pi(X)\right]
  = \E[Y(1)].
\]
The first equality uses consistency ($Y = Y(1)$ when $T{=}1$);
the third uses ignorability ($Y(1) \indep T \mid X$).
\end{proof}

\subsection{The Horvitz--Thompson Estimator}

The corresponding sample estimator is the \emph{Horvitz--Thompson}
(HT) IPW estimator:
\[
  \hat\tau_{\mathrm{IPW}}
  = \frac{1}{n}\sum_{i=1}^n
    \left[\frac{T_i Y_i}{\hat\pi(X_i)}
        - \frac{(1-T_i) Y_i}{1-\hat\pi(X_i)}\right].
\]
This estimator is consistent when $\hat\pi(X)$ converges to $\pi(X)$,
but can be \emph{inefficient} due to extreme weights.  The
\emph{normalised} (Hájek) estimator divides by the sum of weights
within each group and is often more stable in practice:
\[
  \hat\tau_{\mathrm{HJ}}
  = \frac{\sum_i T_i Y_i / \hat\pi(X_i)}
         {\sum_i T_i / \hat\pi(X_i)}
  - \frac{\sum_i (1-T_i) Y_i / (1-\hat\pi(X_i))}
         {\sum_i (1-T_i) / (1-\hat\pi(X_i))}.
\]
The efficient doubly-robust estimator (AIPW) of Chapter~10 combines the
outcome model and the propensity score and achieves the semiparametric
efficiency bound.

% ============================================================
\section{Overlap and Positivity}
% ============================================================

\subsection{The Strict Overlap Condition}

\begin{defn}{Strict Overlap (Positivity)}
\label{w6:defn:overlap}
The \textbf{strict overlap condition} requires:
\[
  0 < \pi(X) < 1 \quad \text{almost surely.}
\]
Every unit has a positive probability of receiving either treatment or
control, regardless of its covariate values.
\end{defn}

\noindent Together with unconfoundedness, strict overlap constitutes
\emph{strong ignorability} (Rosenbaum \& Rubin, 1983).

\medskip\noindent\textbf{Why overlap is necessary.}  The identification
formula
\[
  \tau_{\mathrm{ATE}}
  = \E_X\bigl[\E[Y \mid T{=}1, X] - \E[Y \mid T{=}0, X]\bigr]
\]
requires that both $\E[Y \mid T{=}1, X{=}x]$ and
$\E[Y \mid T{=}0, X{=}x]$ be defined for all $x$ in the support of $X$.
If $\pi(x) = 0$ for some $x$, then no treated units exist with $X = x$,
so $\E[Y \mid T{=}1, X{=}x]$ is not estimable from the data.
Similarly if $\pi(x) = 1$.

\subsection{Practical Consequences of Near-Violations}

Exact violations of overlap ($\pi(X) \in \{0,1\}$) are rare in practice.
More common are \emph{near-violations}: some covariate strata have very
few treated or control units, making the estimated propensity score
close to 0 or 1.

\begin{warn}{Near-Overlap Problems}
When $\hat\pi(X_i)$ is close to 0 or 1, the IPW weight
$1/\hat\pi(X_i)$ or $1/(1 - \hat\pi(X_i))$ is very large.  A small
number of units with extreme weights can dominate the estimator,
increasing variance dramatically.  This is the \emph{practical
positivity violation} problem, distinct from the theoretical violation.
\end{warn}

Near-violations arise when a covariate is a near-perfect predictor of
treatment, when the covariate distributions of treated and control
groups barely overlap, or when the propensity score model is
over-flexible and fits the treatment variable nearly perfectly in-sample.

\subsection{Trimming Strategies}

\subsubsection*{Trimming by Propensity Score}
Restrict the analysis to units with estimated propensity scores in a
prespecified interval $[\eta, 1-\eta]$ for some small $\eta > 0$.
The estimand changes to the ATE in the trimmed subpopulation:
\[
  \tau_{\mathrm{trim}}
  = \E\bigl[Y(1) - Y(0) \mid \eta \le \pi(X) \le 1-\eta\bigr].
\]
The choice of $\eta$ is a bias-variance trade-off: larger $\eta$ reduces
variance but restricts the estimand to a narrower population.

\subsubsection*{Crump et al.\ (2009) Rule}
Crump, Hotz, Imbens \& Mitnik (2009) derive the optimal trimming
threshold that minimizes the asymptotic variance of the trimmed ATE
estimator.  The rule $\eta = 0.1$ --- trimming units with $\pi(X) < 0.1$
or $\pi(X) > 0.9$ --- is a common default.

\subsubsection*{ATT Instead of ATE}
When overlap fails for some treated units but the support of $\pi(X)$
given $T{=}1$ is well-covered by the control group, re-targeting the
estimand to $\tau_{\mathrm{ATT}}$ avoids the problem entirely.  The ATT
requires only $\pi(X) > 0$, not $\pi(X) < 1$.




% ============================================================
\section{Lab: Simulation Study of IPW and Matching Estimators}
\label{w6:sec:lab}
% ============================================================

This lab compares five estimators on a five-covariate data-generating
process with treatment-effect heterogeneity.  With five continuous
confounders, direct stratification is infeasible, making
propensity-score methods the natural choice.  The estimators span two
estimands: the ATE (targeted by the HT and Hájek IPW estimators) and
the ATT (targeted by the ATT-IPW variants and by matching).  Five
questions drive the comparison:
\begin{enumerate}
  \item How much variance does weight normalization save (Hájek over HT)
    for ATE estimation, and why does the baseline outcome mean matter?
  \item Does the same variance problem afflict HT when targeting the ATT,
    and does Hájek normalization help as much?
  \item Among estimators targeting the ATT, do IPW and matching agree
    on the target, and which is more efficient?
  \item How large is the estimand gap between ATT and ATE, and which
    estimators reveal it?
  \item How does replacing the true propensity score with a logistic
    estimate change bias and variance?
\end{enumerate}

% ----------------------------------------
\subsection{Part 1: Correctly Specified Propensity Score}
\label{w6:sec:lab:part1}

\subsubsection*{Data-Generating Process}

\begin{tcolorbox}[defstyle, title={DGP for Lab~6}]
\label{w6:lab:dgp}
Let $X_j \overset{\mathrm{i.i.d.}}{\sim} \mathrm{Normal}(0,1)$ for
$j = 1, \ldots, 5$.  All five covariates enter both the treatment
mechanism and the outcome, so every $X_j$ is a genuine confounder.
The true propensity score is:
\begin{equation}
  \pi(X) = \mathrm{expit}(-0.5 + 0.8X_1 + 0.5X_2 - 0.3X_3 + 0.2X_4 + 0.1X_5),
  \label{w6:lab:ps}
\end{equation}
giving marginal treatment probability $P(T{=}1) \approx 0.40$.  The
potential outcomes are:
\begin{equation}
  Y(t) \;=\; (1 + 0.5X_1)\cdot t
          \;+\; 8 + 0.8X_1 + 0.5X_2 + 0.4X_3 + 0.3X_4 + 0.2X_5 + \varepsilon,
  \qquad \varepsilon \sim \mathrm{N}(0,1).
  \label{w6:lab:po}
\end{equation}
The CATE is $\tau(X) = Y(1) - Y(0) = 1 + 0.5X_1$.  Because $\E[X_1] = 0$:
\begin{equation}
  \tau_{\mathrm{ATE}} = \E[1 + 0.5X_1] = 1.000 \quad (\text{exact}).
  \label{w6:lab:ate}
\end{equation}
Because $\pi(X)$ is increasing in $X_1$, treated units have higher
$X_1$ on average, so $\E[X_1 \mid T{=}1] > 0$ and:
\begin{equation}
  \tau_{\mathrm{ATT}}
  = \E[1 + 0.5X_1 \mid T{=}1]
  \approx 1.199
  \quad (\text{verified by oracle run, } n = 10^7).
  \label{w6:lab:att}
\end{equation}
The gap $\tau_{\mathrm{ATT}} - \tau_{\mathrm{ATE}} \approx 0.199$
reflects a selection effect: units most likely to be treated also
tend to have larger individual treatment effects.  The intercept~$8$
in the outcome model represents a population baseline --- a health
score in a generally healthy population.  Strong ignorability holds
by construction.  The observed outcome is $Y = Y(T)$.
\end{tcolorbox}

With five confounders, direct cell-by-cell stratification is infeasible:
five continuous covariates cannot be coarsened without discarding
information or requiring enormous samples.  The propensity score reduces
this five-dimensional adjustment to a scalar.  The large baseline mean
($\E[Y(0)] = 8$) plays a distinct role: it amplifies the consequences
of extreme IPW weights, making Hájek normalization essential rather
than merely helpful, as the results below demonstrate.

% ----------------------------------------
\subsubsection*{Estimators}

Five estimators are compared, each applied under both the true and
estimated propensity score.  The estimated score $\hat\pi(X)$ is
obtained by logistic regression of $T$ on $(X_1, \ldots, X_5)$, which
is correctly specified since the true logit is linear in all five
covariates.

\medskip\noindent\textbf{Horvitz--Thompson ATE (HT-ATE).}
\[
  \hat\tau_{\mathrm{HT}}
  = \frac{1}{n}\sum_{i=1}^n
    \left[\frac{T_i Y_i}{\hat\pi(X_i)}
        - \frac{(1-T_i)Y_i}{1-\hat\pi(X_i)}\right].
\]
Targets the ATE.  Raw inverse weights can be extreme when $\hat\pi(X_i)$
is near $0$ or $1$, and with $\E[Y(0)]=8$ the variance consequences are
severe.

\medskip\noindent\textbf{Hájek ATE (HJ-ATE).}
\[
  \hat\tau_{\mathrm{HJ}}
  = \frac{\displaystyle\sum_i T_i Y_i / \hat\pi(X_i)}
         {\displaystyle\sum_i T_i   / \hat\pi(X_i)}
  - \frac{\displaystyle\sum_i (1-T_i)Y_i / (1-\hat\pi(X_i))}
         {\displaystyle\sum_i (1-T_i)   / (1-\hat\pi(X_i))}.
\]
Targets the ATE.  Normalizing by the empirical weight sum within each
arm caps the effective influence of any single unit.

\medskip\noindent\textbf{Horvitz--Thompson ATT (HT-ATT).}
The ATT identification formula re-weights the \emph{control} arm to
look like the treated population while leaving the treated arm
unchanged.  The corresponding HT estimator is:
\[
  \hat\tau_{\mathrm{HT,ATT}}
  = \frac{1}{n_1}\sum_{i:\,T_i=1} Y_i
  \;-\;
  \frac{1}{n_1}\sum_{i:\,T_i=0}
    \frac{\hat\pi(X_i)}{1-\hat\pi(X_i)}\,Y_i.
\]
Targets the ATT.  Each control unit receives weight proportional to
its \emph{odds of treatment} $\hat\pi(X_i)/(1-\hat\pi(X_i))$: units
whose covariates resemble the treated population receive large weight.
For a control unit with $\hat\pi(X_i) = 0.95$ the weight is $19$;
combined with a baseline outcome near $8$, this produces extreme
variance --- in practice, even worse than HT-ATE.

\medskip\noindent\textbf{Hájek ATT (HJ-ATT).}
The normalized counterpart replaces the raw-weight sum over the control
arm with a self-normalized version:
\[
  \hat\tau_{\mathrm{HJ,ATT}}
  = \frac{1}{n_1}\sum_{i:\,T_i=1} Y_i
  \;-\;
  \frac{\displaystyle\sum_{i:\,T_i=0}
    \dfrac{\hat\pi(X_i)}{1-\hat\pi(X_i)}\,Y_i}
  {\displaystyle\sum_{i:\,T_i=0}
    \dfrac{\hat\pi(X_i)}{1-\hat\pi(X_i)}}.
\]
Targets the ATT.  The treated arm is already a simple mean and requires
no normalization; the Hájek adjustment applies only to the control arm.
This is the recommended default estimator for the ATT in practice.

\medskip\noindent\textbf{One-to-one nearest-neighbor matching with
replacement (NNM).}
For each treated unit $i$, find the control unit with the closest
propensity score, with reuse allowed \citep{abadie2006large}:
\[
  \hat\jmath(i) = \arg\min_{j:\, T_j=0}
    \bigl|\hat\pi(X_i) - \hat\pi(X_j)\bigr|,
  \qquad
  \hat\tau_{\mathrm{NNM}}
  = \frac{1}{n_1}\sum_{i:\,T_i=1}
    \bigl[Y_i - Y_{\hat\jmath(i)}\bigr].
\]
Targets the ATT.  Matching with replacement ensures every treated unit
finds its globally closest control, preventing the finite-sample
``running out of good matches'' problem that arises without replacement
when the propensity score distribution is wide \citep{abadie2006large}.

% ----------------------------------------
\subsubsection*{Results}

Table~\ref{w6:tab:lab} reports the Monte Carlo mean, bias, standard
deviation, and RMSE across $B = 2{,}000$ replications at $n = 1{,}000$,
seed \texttt{42}.  Bias is computed relative to each estimator's own
target: ATE $= 1.000$ for HT-ATE and HJ-ATE; ATT $\approx 1.199$ for
HT-ATT, HJ-ATT, and NNM.

\begin{table}[h]
\centering
\renewcommand{\arraystretch}{1.5}
\caption{Monte Carlo performance of five propensity-score estimators
  under known and estimated propensity scores.
  True ATE $= 1.000$ (exact); true ATT $\approx 1.199$ (oracle, $n=10^7$).
  $n = 1{,}000$, $B = 2{,}000$ replications, seed \texttt{42}.}
\label{w6:tab:lab}
\begin{tabular}{llccccc}
\toprule
\textbf{PS} & \textbf{Estimator} & \textbf{Estimand}
  & \textbf{Mean} & \textbf{Bias} & \textbf{SD} & \textbf{RMSE} \\
\midrule
Known
  & HT-ATE  & ATE & 1.008 & $+$0.008 & 0.626 & 0.626 \\
  & HJ-ATE  & ATE & 1.003 & $+$0.003 & 0.145 & 0.145 \\[3pt]
  & HT-ATT  & ATT & 1.184 & $-$0.016 & 0.725 & 0.725 \\
  & HJ-ATT  & ATT & 1.203 & $+$0.003 & 0.131 & 0.131 \\
  & NNM     & ATT & 1.207 & $+$0.008 & 0.130 & 0.130 \\
\midrule
Estimated
  & HT-ATE  & ATE & 0.998 & $-$0.002 & 0.194 & 0.194 \\
  & HJ-ATE  & ATE & 1.002 & $+$0.002 & 0.102 & 0.103 \\[3pt]
  & HT-ATT  & ATT & 1.202 & $+$0.002 & 0.251 & 0.251 \\
  & HJ-ATT  & ATT & 1.202 & $+$0.002 & 0.101 & 0.101 \\
  & NNM     & ATT & 1.205 & $+$0.006 & 0.121 & 0.121 \\
\bottomrule
\end{tabular}
\end{table}

\medskip
Five lessons stand out.

\medskip
\noindent\textbf{1.\ Hájek normalization is essential for ATE
estimation when outcomes are non-centered.}
With the known propensity score, HT-ATE achieves SD $= 0.626$, while
HJ-ATE achieves SD $= 0.145$ --- a \textbf{4.3-fold reduction} from
simply normalizing the weights.  The mechanism is transparent: a unit
with $\pi(X_i) = 0.05$ receives raw weight $1/0.05 = 20$; multiplied
by an outcome near $\E[Y(0)] = 8$, its contribution to the HT sum is
of order $160$.  Variability in which units are extreme across samples
inflates the HT standard deviation to more than half the true effect
size.  The Hájek estimator caps each unit's effective weight at its
share of the total within its arm, suppressing this mechanism.

\medskip
\noindent\textbf{2.\ HT-ATT is even more unstable; Hájek normalization
is equally essential for ATT estimation.}
With the known PS, HT-ATT achieves SD $= 0.725$, slightly \emph{worse}
than HT-ATE (SD $= 0.626$).  The reason is structural: the ATT
re-weights the control arm by the treatment odds $\pi(X)/(1-\pi(X))$.
A control unit with $\pi(X) = 0.95$ receives an odds-ratio weight of
$0.95/0.05 = 19$; combined with the large baseline outcome this is as
explosive as the ATE inverse-probability weight.  The apparent bias of
$-0.016$ for HT-ATT is within one Monte Carlo standard error of
$0.725/\sqrt{2000} \approx 0.016$ and should be attributed to
simulation noise rather than systematic inconsistency.  HJ-ATT reduces
SD to $0.131$ --- a \textbf{5.5-fold reduction} --- confirming that
Hájek normalization is not an ATE-specific fix but a general principle
applicable whenever inverse-probability-style weights are used.

\medskip
\noindent\textbf{3.\ HJ-ATT and NNM converge to the same target with
nearly identical efficiency.}
Under the known PS, HJ-ATT and NNM both target the ATT $\approx 1.199$
and achieve essentially the same standard deviation (0.131 vs.\ 0.130).
This is the sharpest demonstration of the ``two routes, one estimand''
principle: IPW re-weights the full sample using the propensity score;
matching discards some control units and retains only those nearest
each treated unit on the propensity scale.  Both are consistent for
the ATT under strong ignorability, and at $n = 1{,}000$ they are
virtually indistinguishable in efficiency.  The practical implication
is that the choice between IPW-ATT and matching can be made on
secondary grounds --- interpretability, balance diagnostics, or
computational cost --- rather than efficiency when the PS is known.

\medskip
\noindent\textbf{4.\ The estimand gap between ATE and ATT is large and
clearly revealed by the table.}
The ATE estimators (HT-ATE, HJ-ATE) converge to $1.000$; the ATT
estimators (HT-ATT, HJ-ATT, NNM) converge to $\approx 1.199$.  The
$0.199$ gap is not a bias --- it is a real difference in the
estimands.  High-$X_1$ units are more likely to be treated
(equation~\eqref{w6:lab:ps}) and have larger treatment effects
(equation~\eqref{w6:lab:po}); the treated subpopulation therefore
benefits more than the full population average.  A researcher who
applies NNM and reports its output against the ATE benchmark of $1.000$
would conclude the estimator has $20\%$ bias; in fact it is unbiased
for the correct target.  Declaring the estimand before looking at any
results is a non-negotiable discipline.

\medskip
\noindent\textbf{5.\ Estimated PS collapses HT variance and makes
HJ-ATT the most efficient ATT estimator.}
Under the estimated PS, HT-ATE SD falls from $0.626$ to $0.194$ (69\%
reduction); HT-ATT SD falls from $0.725$ to $0.251$ (65\% reduction).
In both cases the shrinkage of fitted logistic probabilities toward the
sample mean of $T$ trims the most extreme weights automatically.  More
strikingly, HJ-ATT with estimated PS achieves SD $= 0.101$ --- beating
both NNM (SD $= 0.121$) and HT-ATT (SD $= 0.251$) by a wide margin and
matching HJ-ATE (SD $= 0.102$) almost exactly.  This is the finite-sample
analog of the efficiency result of \citet{imbens2015causal}: a
properly normalized IPW estimator for the ATT is not only unbiased but
can be more efficient than matching when the propensity score model is
well-estimated.

\begin{warn}{The Estimand Must Be Chosen Before the Estimator}
In this DGP, $\tau_{\mathrm{ATE}} = 1.000$ and
$\tau_{\mathrm{ATT}} \approx 1.199$.  The gap arises because treatment
status is correlated with the individual treatment effect: high-$X_1$
units are simultaneously more likely to be treated
(equation~\eqref{w6:lab:ps}) and more likely to benefit
(equation~\eqref{w6:lab:po}).  This situation --- treatment-effect
heterogeneity correlated with propensity for treatment
\citep{imbens2015causal} --- is common in practice: individuals often
select into treatment partly because they anticipate larger benefits.

IPW-ATE and HJ-ATE re-weight the full sample to represent the
population; IPW-ATT and HJ-ATT re-weight only the control arm to
represent the treated population; NNM directly pairs each treated unit
with a similar control.  All three ATT estimators target the same
quantity via different mechanisms and deliver consistent results here.
Choosing among them should be driven by the policy question (is the
ATE or ATT the right estimand?) and by practical concerns about balance
and model specification, not by which produces the larger or smaller
point estimate.
\end{warn}

\begin{remark}[What Part~1 Does Not Show]
\label{w6:rem:lab-limits}
Part~1 uses a correctly specified logistic propensity model, which is
the best-case scenario for IPW.  In practice the true PS may be
nonlinear, making correct logistic specification impossible with the
predictors at hand.  Part~2 below addresses exactly this setting.  The
doubly-robust AIPW estimator of Chapter~10 provides a complementary
fix: it is consistent if \emph{either} the PS model or the outcome
model is correctly specified, offering robustness that neither pure IPW
nor pure matching can achieve alone.
\end{remark}

% ----------------------------------------
\subsection{Part 2: Robustness to PS Model Misspecification}
\label{w6:sec:lab:part2}

Part~1 established that under a correctly specified PS model, IPW and
NNM are both consistent for their respective estimands, with HJ-ATT
and NNM delivering similar efficiency.  Part~2 breaks that assumption
by introducing a \emph{nonlinear} true propensity score that the
standard logistic regression cannot capture.

\medskip\noindent\textbf{Modified DGP.}
The outcome model is unchanged from equation~\eqref{w6:lab:po}.  The
true propensity score now contains a quadratic term in $X_1$, and $X_2$
and $X_3$ do not affect treatment:
\begin{equation}
  \pi^*(X) = \mathrm{expit}\!\bigl(-0.5 + 0.8X_1 + 0.5X_1^2 + 0.2X_4 + 0.1X_5\bigr).
  \label{w6:lab:ps2}
\end{equation}
The quadratic $0.5X_1^2$ means that both large positive \emph{and}
large negative values of $X_1$ increase the probability of treatment,
creating a U-shaped propensity surface.  Because $\E[X_1] = 0$ and
the CATE is still $\tau(X) = 1 + 0.5X_1$:
\[
  \tau_{\mathrm{ATE}} = 1.000 \quad (\text{exact, unchanged}).
\]
The ATT changes because the new PS shifts which units are treated:
\begin{equation}
  \tau_{\mathrm{ATT}}^* \approx 1.152
  \quad (\text{oracle run, }n = 10^7).
  \label{w6:lab:att2}
\end{equation}
The \emph{estimated} PS is still obtained by fitting a linear logistic
regression of $T$ on $(X_1, X_2, X_3, X_4, X_5)$.  This model is
\textbf{misspecified}: it omits the $X_1^2$ term and includes $X_2$
and $X_3$, which do not appear in the true logit.

\medskip\noindent\textbf{Results.}
Table~\ref{w6:tab:lab2} reports the same five estimators as Part~1,
now applied under the modified DGP.  The ``True PS'' rows use
$\pi^*(X)$ from equation~\eqref{w6:lab:ps2}; the ``Misspecified''
rows use a linear logistic estimate $\hat\pi(X)$ that omits $X_1^2$.

\begin{table}[h]
\centering
\renewcommand{\arraystretch}{1.5}
\caption{Part~2: Monte Carlo performance under a nonlinear true PS and
  a misspecified linear logistic estimate.  True ATE $= 1.000$ (exact);
  true ATT $\approx 1.152$ (oracle, $n=10^7$, new DGP).
  $n = 1{,}000$, $B = 2{,}000$ replications, seed \texttt{42}.}
\label{w6:tab:lab2}
\begin{tabular}{llccccc}
\toprule
\textbf{PS} & \textbf{Estimator} & \textbf{Estimand}
  & \textbf{Mean} & \textbf{Bias} & \textbf{SD} & \textbf{RMSE} \\
\midrule
True
  & HT-ATE  & ATE & 1.028 & $+$0.028 & 0.824 & 0.825 \\
  & HJ-ATE  & ATE & 1.013 & $+$0.013 & 0.152 & 0.152 \\[3pt]
  & HT-ATT  & ATT & 1.186 & $+$0.034 & 1.333 & 1.333 \\
  & HJ-ATT  & ATT & 1.183 & $+$0.031 & 0.212 & 0.215 \\
  & NNM     & ATT & 1.178 & $+$0.026 & 0.175 & 0.177 \\
\midrule
Misspecified
  & HT-ATE  & ATE & 1.478 & $+$0.478 & 0.147 & 0.501 \\
  & HJ-ATE  & ATE & 0.861 & $-$0.139 & 0.087 & 0.164 \\[3pt]
  & HT-ATT  & ATT & 1.780 & $+$0.628 & 0.162 & 0.649 \\
  & HJ-ATT  & ATT & 1.306 & $+$0.154 & 0.081 & 0.174 \\
  & NNM     & ATT & 1.172 & $+$0.020 & 0.138 & 0.139 \\
\bottomrule
\end{tabular}
\end{table}

\medskip
Three lessons complete the picture.

\medskip
\noindent\textbf{6.\ PS misspecification makes all IPW estimators
inconsistent; the Hájek correction cannot fix this.}
Under the misspecified PS, every IPW estimator is substantially biased.
HT-ATE bias is $+0.478$, nearly half the true effect; HJ-ATE bias is
$-0.139$.  For the ATT estimators the problem is even worse: HT-ATT
bias is $+0.628$; HJ-ATT bias is $+0.154$.  Comparing these to the
near-zero biases in the ``True PS'' rows makes clear that the Hájek
normalization is a \emph{variance} correction --- it cannot compensate
for a misspecified model that produces systematically wrong weights.
IPW consistency is tied to consistency of the PS estimator; a
structurally wrong model produces structurally wrong weights, and no
normalization of wrong weights produces correct estimates.

Notice also that HJ-ATE and HT-ATE have biases in \emph{opposite
directions} ($-0.140$ vs $+0.481$).  The true PS is U-shaped in $X_1$:
extreme values of $X_1$ in either direction attract high treatment
probability.  The misspecified linear logistic underestimates $\pi^*$
for large $|X_1|$ and overestimates it for moderate $X_1$.  HT-ATE
up-weights the treated arm by $1/\hat\pi$ and down-weights the control
arm by $1/(1-\hat\pi)$; HJ-ATE applies normalized versions of the same
weights.  The two estimators inherit the misfit in opposite ways
because they respond differently to the systematic over/under-estimation
of the score.

\medskip
\noindent\textbf{7.\ NNM is substantially more robust to PS
misspecification.}
The NNM bias barely changes: $+0.026$ under the true PS versus $+0.020$
under the misspecified PS.  Its RMSE actually \emph{improves} slightly
($0.177 \to 0.139$) because the misspecified logistic scores, being
shrunk toward the mean of $T$, produce somewhat less variable matches
than the true nonlinear scores do.  This robustness has a direct
explanation: NNM does not require the PS model to be correct in order
to produce good matches --- it only requires that the estimated score
roughly \emph{orders} units so that matched pairs are approximately
balanced on the true confounders.  The linear logistic model, though
wrong, still captures the dominant linear effects of $X_1$, $X_4$, and
$X_5$ on treatment; the quadratic component is missed, but the ranking
of units is not drastically distorted.  As a result, matched pairs
remain reasonably balanced on the variables that enter the outcome, and
the bias stays small.
This robustness of matching estimators to propensity score model
misspecification is studied theoretically in
\citet{yang2016multilevel} and \citet{yang2023multiply}.

\medskip
\noindent\textbf{8.\ The relative advantage of matching over IPW
reverses with model misspecification.}
In Part~1 (correctly specified PS), HJ-ATT was more efficient than NNM
under the estimated PS (SD $= 0.101$ vs $0.121$).  In Part~2
(misspecified PS), NNM dominates on both bias \emph{and} RMSE
($0.020$, $0.139$) relative to HJ-ATT ($0.154$, $0.174$).  The
comparison is summarized in Table~\ref{w6:tab:comparison}.  The
practical implication is that the efficiency advantage of IPW is
contingent on the PS model being approximately correct.  When there is
reason to doubt the functional form of the propensity model --- for
instance, when treatment is driven by a continuous covariate in a
nonlinear way --- matching provides a more robust alternative, because
it only relies on local covariate similarity rather than a globally
correct weighting function.

\begin{table}[h]
\centering
\renewcommand{\arraystretch}{1.4}
\caption{Head-to-head comparison of HJ-ATT and NNM under correct
  vs.\ misspecified PS estimation (estimated PS column from each part).}
\label{w6:tab:comparison}
\begin{tabular}{lcccc}
\toprule
& \multicolumn{2}{c}{\textbf{Part~1: Correct PS model}}
& \multicolumn{2}{c}{\textbf{Part~2: Misspecified PS model}} \\
\cmidrule(lr){2-3}\cmidrule(lr){4-5}
\textbf{Estimator} & Bias & RMSE & Bias & RMSE \\
\midrule
HJ-ATT & $+0.002$ & $0.101$ & $+0.154$ & $0.174$ \\
NNM    & $+0.006$ & $0.121$ & $+0.020$ & $0.139$ \\
\bottomrule
\end{tabular}
\end{table}

\begin{warn}{IPW Needs a Correct PS Model; Matching Needs Covariate Balance}
The two lessons of this lab can be compressed into a single
contrast.  IPW is a \emph{model-dependent} estimator: its consistency
relies on the PS model being correctly specified.  Get the model right
and IPW (especially the Hájek form) is highly efficient; get it wrong
and the bias from incorrect weights can be large and systematic,
as the $+0.154$ HJ-ATT bias in Part~2 demonstrates.

Matching is a \emph{design-dependent} estimator: its consistency relies
on matched pairs being approximately balanced on the true confounders.
This condition can hold even when the PS model is wrong, provided the
estimated score still separates units well enough that close matches on
the score imply closeness on the full covariate vector.  The cost is
efficiency: when the PS model is correct, matching leaves variance on
the table relative to IPW.

Neither dominates unconditionally.  The practical recommendation is to
assess the plausibility of the PS model (e.g., via balance diagnostics
after stratification or matching) before choosing between IPW and
matching as the primary estimator --- or to use both as a sensitivity
check \citep{yang2016multilevel, yang2023multiply}.
\end{warn}

\section{The Limits of Propensity Score Methods}
% ============================================================

\subsection{The Untestable Assumption}

Every result in this chapter --- the Propensity Score Theorem, the IPW
identification formula, the matching estimator --- rests on the
unconfoundedness assumption:
\[
  \bigl(Y(0),\, Y(1)\bigr) \;\indep\; T \;\Big|\; X.
\]
This says that \emph{all} confounders of the $T$--$Y$ relationship are
observed and included in $X$.  In graphical terms: $X$ blocks every
back-door path from $T$ to $Y$.

\begin{tcolorbox}[defstyle, title={Design-Based vs.\ Model-Based Identification}]
Randomization \emph{guarantees} ignorability: by construction, no
back-door paths exist, so the interventional and observational
distributions coincide.  Propensity scores \emph{assume} ignorability:
the analyst hopes that all confounders have been measured and included
in $X$, but this can never be verified from the data alone.  This
distinction --- between identification secured by the study design and
identification secured by a modeling assumption --- is one of the most
important dividing lines in causal inference.
\end{tcolorbox}

\subsection{The Road to Instrumental Variables}

The limitation of back-door adjustment motivates the search for
alternative identification strategies when unobserved confounders are
present.  The canonical solution is the instrumental variable (IV),
studied in Chapter~7.

\medskip
\renewcommand{\arraystretch}{1.5}
\begin{tabular*}{\linewidth}{@{\extracolsep{\fill}} p{4.2cm} p{5.0cm} p{5.4cm}}
  \toprule
  \textbf{Strategy}
    & \textbf{Key assumption}
    & \textbf{Identification mechanism} \\
  \midrule
  Back-door adjustment\newline (propensity score)
    & All confounders are observed
    & Condition on $X$ to block $T \leftarrow U \rightarrow Y$ \\[4pt]
  Instrumental variables
    & Some confounders may be unobserved
    & Exploit exogenous variation $Z \rightarrow T$ \\
  \bottomrule
\end{tabular*}
\renewcommand{\arraystretch}{1}

\medskip\noindent Back-door adjustment removes confounding by
conditioning on observed variables.  Instrumental variables take a
different approach: instead of blocking the confounding path
$T \leftarrow U \rightarrow Y$, they identify the causal effect using
variation in $T$ induced by an external source $Z$ that is independent
of $U$.  Concretely, the ratio of the $Z$-induced change in $Y$ to the
$Z$-induced change in $T$ recovers the causal effect --- this ratio is
the Wald estimator, derived formally in Chapter~7.

\medskip
\begin{center}
\begin{tikzpicture}[
  node distance=16mm and 22mm,
  ivbox/.style={rectangle, rounded corners=3pt, minimum width=32mm,
                minimum height=10mm, align=center, font=\small,
                draw=accent, line width=0.8pt},
  ivlbl/.style={font=\footnotesize\itshape\color{isublue}},
  ivarr/.style={-{Stealth[length=5pt]}, line width=0.9pt, color=accent},
  ivredarr/.style={-{Stealth[length=5pt]}, line width=0.9pt,
                   color=backdoorred, dashed}
]
  \node[ivbox, fill=defbg]                (Z)  {Instrument $Z$};
  \node[ivbox, fill=defbg,  right=of Z]  (T)  {Treatment $T$};
  \node[ivbox, fill=defbg,  right=of T]  (Y)  {Outcome $Y$};
  \node[ivbox, fill=warnbg, above=10mm of T, xshift=15mm]
    (U) {Unobserved\\[-2pt]Confounder $U$};

  \draw[ivarr]    (Z) -- node[below, ivlbl]{Relevance} (T);
  \draw[ivarr]    (T) -- (Y);
  \draw[ivredarr] (U) -- (T);
  \draw[ivredarr] (U) -- (Y);

  \node[ivlbl, below=4mm of Z, align=center]
    {\emph{Exogenous}: $Z \indep U$\\
     Exclusion: $Z$ affects $Y$\\only through $T$};
\end{tikzpicture}
\end{center}
\medskip

% ============================================================
\section{Summary}
% ============================================================

\begin{enumerate}
  \item \textbf{The propensity score reduces dimension.}  Under
    unconfoundedness, $\pi(X) = P(T{=}1 \mid X)$ is a \emph{balancing
    score}: $T \indep X \mid \pi(X)$.  Consequently,
    $(Y(0),Y(1)) \indep T \mid \pi(X)$ (Theorem~\ref{w6:thm:ps}), so
    identification requires adjustment for the scalar $\pi(X)$ alone, not
    the full vector $X$.

  \item \textbf{IPW identification.}  The ATE is identified as
    $\E[TY/\pi(X)] - \E[(1-T)Y/(1-\pi(X))]$ under strong ignorability.
    The Horvitz--Thompson IPW estimator is consistent but can be
    sensitive to extreme weights; the Hájek and doubly-robust (AIPW)
    variants are more efficient.

  \item \textbf{Matching targets the ATT.}  Propensity score matching
    finds control units with similar propensity scores to treated units.
    It targets the ATT and provides transparent covariate balance
    diagnostics, but differs fundamentally from IPW weighting in
    estimand and methodology.

  \item \textbf{Strict overlap is non-negotiable.}  The condition
    $0 < \pi(X) < 1$ a.s.\ is necessary for identification of the ATE.
    Near-violations create extreme IPW weights and large variance;
    trimming strategies address this at the cost of restricting the
    estimand.

  \item \textbf{The fundamental limitation.}  Unconfoundedness
    $(Y(0),Y(1)) \indep T \mid X$ is a strong, untestable assumption.
    When unobserved confounders are present, propensity score adjustment
    fails and an instrumental variable is needed to identify the causal
    effect.
\end{enumerate}

\begin{tcolorbox}[defstyle, title={Comparison of Identification Strategies}]
\begin{center}
\renewcommand{\arraystretch}{1.4}
\begin{tabular}{p{3.5cm} p{4.5cm} p{4.5cm}}
\toprule
\textbf{Design} & \textbf{Key assumption} & \textbf{Identified estimand} \\
\midrule
Randomized experiment
  & $(Y(0),Y(1)) \indep T$\newline (by design)
  & ATE \\[4pt]
Propensity-score adjustment
  & $(Y(0),Y(1)) \indep T \mid X$\newline (unconfoundedness)
  & ATE or ATT \\[4pt]
Instrumental variables
  & Relevance, exogeneity, exclusion
  & LATE (compliers) \\
\bottomrule
\end{tabular}
\end{center}
\renewcommand{\arraystretch}{1}
\end{tcolorbox}

% ============================================================
\section*{Problems}
\addcontentsline{toc}{section}{Problems}
% ============================================================

\begin{enumerate}

\item \textbf{Propensity score and balancing.}
  Suppose $X = (X_1, X_2)$, where $X_1$ is binary and $X_2$ is
  continuous.  The propensity score model is correctly specified as
  $\mathrm{logit}(\pi(X)) = \beta_0 + \beta_1 X_1 + \beta_2 X_2$.
  \begin{enumerate}
    \item State and prove the Balancing Property $T \indep X \mid \pi(X)$.
    \item Two units $i$ and $j$ have $X_i = (1, 2.3)$ and
      $X_j = (0, 3.5)$ but $\hat\pi(X_i) = \hat\pi(X_j) = 0.4$.
      If $i$ is treated and $j$ is control, explain why comparing $Y_i$
      and $Y_j$ is a valid approximation to the counterfactual
      comparison, and what assumption is required.
    \item Explain why matching on $\hat\pi(X)$ is \emph{not} the same
      as matching on $X$ directly.  Under what conditions do they give
      the same answer?
  \end{enumerate}

\item \textbf{ATE, ATT, and propensity score weighting.}
  A dataset has $n = 1000$ observations.  The estimated propensity
  scores are $\hat\pi(X_i)$ for $i = 1, \ldots, n$, with $n_1 = 400$
  treated units.
  \begin{enumerate}
    \item Write the Horvitz--Thompson IPW estimator of the ATE as a
      weighted sum of observed outcomes $Y_i$.  What weights do treated
      units receive?  What weights do control units receive?
    \item Derive an analogous IPW estimator for the ATT.
      [\emph{Hint}: the ATT averages over the treated distribution; the
      denominator of the weight for control units should reflect this.]
    \item Show that when the propensity score is constant
      ($\pi(X) = p$ for all units), the IPW estimator of the ATE reduces
      to the difference-in-means estimator.  Under what experimental
      design does $\pi(X) = 0.5$ exactly?
  \end{enumerate}

\item \textbf{Overlap.}
  Let $\pi(X)$ be the true propensity score and suppose overlap fails at
  $X = x_0$: $\pi(x_0) = 1$.
  \begin{enumerate}
    \item Which causal quantity ($\tau_{\mathrm{ATE}}$ or
      $\tau_{\mathrm{ATT}}$) cannot be identified?  Give the exact
      formula that breaks down and explain why.
    \item Suppose overlap fails only on a set $\mathcal{S}$ of covariate
      values with $P(X \in \mathcal{S}) = 0.15$.  Define the
      \emph{trimmed ATE} estimand.  How does it differ from the ATE?
    \item A researcher applies the Crump et al.\ (2009) rule with
      $\eta = 0.1$, removing all units with $\hat\pi(X_i) < 0.1$ or
      $\hat\pi(X_i) > 0.9$.  This removes 8\% of the sample.
      \begin{enumerate}
        \item[(i)] List two reasons the trimmed estimator may have
          lower variance than the untrimmed IPW estimator.
        \item[(ii)] What is the cost in terms of external validity?
      \end{enumerate}
  \end{enumerate}

\item \textbf{Sensitivity analysis (conceptual).}
  Suppose you estimate $\hat\tau_{\mathrm{ATE}} = 3.2$ using
  propensity-score methods, and a referee questions whether
  unconfoundedness holds.
  \begin{enumerate}
    \item Define what it means for $U$ to be a ``hidden confounder'' in
      the context of the DAG, and explain why its presence invalidates
      the IPW identification formula.
    \item Rosenbaum (2002) proposes quantifying the sensitivity to an
      unobserved binary confounder $U$ through a sensitivity parameter
      $\Gamma \ge 1$, where $\Gamma = 1$ corresponds to no hidden bias.
      Explain in words what $\Gamma = 2$ means for the odds of
      treatment.
    \item Without doing formal calculations, describe the narrative a
      researcher might construct to argue that their estimated effect of
      3.2 is \emph{robust} to a confounder with $\Gamma \le 1.5$.
    \item Why does the IV strategy of Chapter~7 avoid this sensitivity
      problem entirely, and what assumption replaces unconfoundedness?
  \end{enumerate}

\item \textbf{Matching vs.\ IPW.}
  You have $n = 500$ observations with a binary treatment and a
  scalar propensity score.  You compare two estimators:
  (A) Hájek IPW estimator targeting the ATE, and
  (B) 1:1 nearest-neighbor matching targeting the ATT.
  \begin{enumerate}
    \item Explain in one sentence why estimator (A) and estimator (B)
      target different estimands even though both use $\hat\pi(X)$.
    \item In your matched sample, the standardized mean difference (SMD)
      for covariate $X_1$ is 0.05 after matching and 0.45 before
      matching.  What does this tell you about the success of matching,
      and what assumption does balance on observed covariates not verify?
    \item Under what condition on the treatment effect heterogeneity
      are the ATE and ATT equal?
  \end{enumerate}

\end{enumerate}

% Bibliography (standalone only)
\ifSubfilesClassLoaded{%
  \bibliographystyle{plainnat}
  \bibliography{causal_inference}
}{}

\end{document}